\documentclass[a4paper,11pt]{article}
\usepackage[T1]{fontenc}
\usepackage[utf8]{inputenc}
\usepackage{t1enc}
\usepackage[dvips]{color}
\usepackage{amsmath,amssymb,makeidx,stmaryrd}
\usepackage[normalem]{ulem}
\usepackage{fancybox}
\definecolor{gris}{gray}{0.9}
\usepackage{pst-plot,pst-tree,pstricks,pst-node,pst-text}
\usepackage{pst-eucl}
\usepackage{pstricks-add}
\usepackage[left=3.5cm, right=3.5cm, top=3cm, bottom=3cm]{geometry}
\newcommand{\vect}[1]{\overrightarrow{\,\mathstrut#1\,}}
\def\bas{\fcolorbox{gris}{gris}{\parbox{0.986\textwidth}{
	\color{red}\footnotesize 1992\hfill Page : \thepage/1}}}
\def\haut{\fcolorbox{gris}{gris}{\parbox{0.986\textwidth}{
	\color{red}
	\bf Concours Général de Mathématiques \hfill 1992
	}}}
\makeatletter
\def\ps@cg{%
\let\@oddhead\haut%
\let\@evenhead\haut%
\let\@oddfoot\bas
\let\@evenfoot\bas}
\makeatother
\newcommand{\R}{\mathbb{R}}
\newcommand{\N}{\mathbb{N}}
\newcommand{\D}{\mathbb{D}}
\newcommand{\Z}{\mathbb{Z}}
\newcommand{\Q}{\mathbb{Q}}
\newcommand{\C}{\mathbb{C}}
\pagestyle{cg}
\renewcommand{\theenumi}{\textbf{\arabic{enumi}}}
\renewcommand{\labelenumi}{\textbf{\theenumi.}}
\renewcommand{\theenumii}{\textbf{\alph{enumii}}}
\renewcommand{\labelenumii}{\textbf{\theenumii.}}
\parindent0pt
\everymath{\displaystyle}
\usepackage[french]{babel}

\DecimalMathComma
\usepackage[np]{numprint}
\begin{document}

\begin{center}\textbf{{\color{blue} Exercice I}}
\end{center}

On dit qu'une partie non vide $\Delta$ du plan $\mathcal{P}$ est convexe si elle possède la propriété suivante :

\begin{center}
Pour tout couple (M, N) de points de $\Delta$, le segment [MN] est contenu dans $\Delta$
\end{center}

Soit $\Delta$ une partie convexe du plan $\mathcal{P}$. Si A est un point de $\mathcal{P}$, à tout couple (M, N) de points de $\Delta$ on associe le point $m$ défini par

\[\overline{\text{A}m} = \dfrac12 \overline{\text{MN}}\]

et on note $\delta_{\text{A}}(\Delta)$ l'ensemble des points $m$ ainsi obtenus.

\medskip

\begin{enumerate}
\item 
	\begin{enumerate}
		\item Montrer que $\delta_{\text{A}}(\Delta)$ admet un centre de symétrie.
		\item À quelle condition a-t-on $\delta_{\text{A}}(\Delta) =\Delta$ ?
		\item Soient B et C deux points du plan $\mathcal{P}$. Par quelle transformation passe-t-on de $\delta_{\text{B}}(\Delta)$ à $\delta_{\text{C}}(\Delta)$ ?
	\end{enumerate}	
\item Soit A un point du plan $\mathcal{P}$. Déterminer et représenter $\delta_{\text{A}}(\Delta)$ lorsque:
	\begin{enumerate}
		\item $\Delta$ est une bande du plan $\mathcal{P}$ délimitée par deux droites parallèles.
		\item $\Delta$ est délimitée par un triangle.
		\item $\Delta$ est un demi-disque.
	\end{enumerate}	
\item Montrer que dans les deux cas 2. b. et 2. c., les contours de $\Delta$ et $\delta_{\text{A}}(\Delta)$ ont la même longueur.
\end{enumerate}

\begin{center}\textbf{{\color{blue} Exercice II}}
\end{center}

Soit $(C)$ un cercle du plan de rayon l.

\medskip

\begin{enumerate}
\item Déterminer les triangles ABC inscrits dans le cercle $(C)$ pour lesquels la somme

\begin{center}AB$^2$ + BC$^2$ + CA$^2$ \end{center}
est maximale.
\item Déterminer les quadrilatères ABCD inscrits dans le cercle $(C)$ pour lesquels la somme

\begin{center}AB$^2$ + AC$^2$ + AD$^2$ + BC$^2$ + BD$^2$ + CD$^2$ \end{center}

est maximale. Représenter un tel quadrilatère.
\end{enumerate}

\begin{center}\textbf{{\color{blue} Exercice III}}
\end{center}


Soit ABCD un tétraèdre inscrit dans une sphère de centre O. 

On note G l'isobarycentre des quatre sommets du tétraèdre et I le centre de la sphère inscrite dans le tétraèdre. Montrer que les trois propriétés suivantes sont équivalentes :

\medskip

\begin{enumerate}
\item Les deux points O et G sont confondus.
\item Les quatre faces du tétraèdre sont isométriques. 
\item Les deux points O et I sont confondus.
\end{enumerate}

\begin{center}\textbf{{\color{blue} Exercice IV}}
\end{center}


Soit $\left(u_n\right)$ la suite numérique définie par la donnée de ses deux premiers termes $u_0$ et $u_1,\: 0 < u_0 < 1$ et $0 < u_1 < 1$, et la relation de récurrence:

\[u_{n+2} = \dfrac12 \left(\sqrt{u_{n+1}} + \sqrt{u_n}\right).\]

\begin{enumerate}
\item Montrer que la suite $\left(u_n\right)$ est convergente et déterminer sa limite.
\item Montrer qu'à partir d'un certain rang $n_0$, la suite $\left(u_n\right)$ est monotone ( on ne demande pas de déterminer $n_0$ qui dépend des valeurs initiales $u_0$ et $u_1$).
\end{enumerate}

\begin{center}\textbf{{\color{blue} Exercice V}}
\end{center}

Quel est le chiffre des unités du plus grand nombre entier inférieur ou égal à 
$\dfrac{10^{\np{1992}}}{10^{83} + 7}$ ?
\end{document}